\section{Design Choices and Information-Flow Argument}

\subsection{Design Scope and Priority}

The project prioritizes formal precision over feature breadth. We therefore choose a compact model with:
\begin{itemize}
    \item one fixed actor universe $\mathcal{A}$,
    \item decentralized labels on all persistent objects,
    \item white-box user programs only,
    \item a small trusted set of server-side declassifiers.
\end{itemize}

This avoids unverifiable corner cases from arbitrary black-box user code.

\subsection{Why a Decentralized Label Model}

Multiple healthcare stakeholders have different security interests: patients, doctors, hospitals, researchers, and regulators. Decentralized labels allow each owner to impose independent constraints without requiring global authority coordination. For example, patient data might have labels $(patient:\{doctor\})$ (patient's read policy), $(hospital:\{hospital,researcher\})$ (hospital's policy), and $(regulator:\{hospital\})$ (regulatory policy). A reader must satisfy all three. This is more expressive and compositional than centralized approaches.

\subsection{Flow Order and Program Rules}

We write $L_1 \sqsubseteq L_2$ when information is allowed to flow from label $L_1$ to $L_2$, respecting both confidentiality and integrity components. This partial order is the foundation of all flow safety checks.

Information flows upward in the order: if data has label $L_1$ and the destination has label $L_2$ with $L_1 \sqsubseteq L_2$, then copying or deriving data at $L_1$ into the destination is safe. The check prevents both explicit flows (direct assignments like $x \leftarrow y$) and implicit flows (dependencies through control flow).

Core program rules used by the server checker:
\[
\frac{\Gamma \vdash e : L_e \qquad pc \sqcup L_e \sqsubseteq \Gamma(x)}{\Gamma,pc \vdash x := e}
\]

where $\Gamma$ is a type environment mapping variables to labels, $pc$ is the ``program counter label'' representing the security level of the current control context, $L_e$ is the label inferred for expression $e$, and $\sqcup$ is the join operation (combining two labels into the least upper bound).

For conditionals:
\[
\frac{\Gamma, pc \vdash \text{cond} : L_{cond} \quad \Gamma, (pc \sqcup L_{cond}) \vdash B}{\Gamma, pc \vdash \text{if cond then } B}
\]

The $pc$ is raised within the branch to $pc \sqcup L_{cond}$, which accounts for the implicit dependency on the condition. This blocks implicit information leaks where an attacker observes which branch was taken by watching which variables are assigned.

\textbf{Example: Blocking an Implicit Leak}. Suppose we have:
\begin{itemize}
    \item $secret: (confidential:\{doctor\})$
    \item $public: (confidential:\{\})$ (readable by all)
\end{itemize}

If a user writes:
\[\text{if } (secret > 0) \text{ then } public := 0 \text{ else } public := 1\]

The type checker would compute:
\begin{enumerate}
    \item Condition type: $L_{cond} = (confidential:\{doctor\})$ from $secret$.
    \item Program counter inside branch: $pc \sqcup L_{cond} = (confidential:\{doctor\})$.
    \item Assignment check: Does $(confidential:\{doctor\}) \sqsubseteq (confidential:\{\})$? No, because $\{doctor\} \not\supseteq \{\}$ (the reader set is not a superset). \times
\end{enumerate}

The assignment is rejected at check time, preventing the leak.

\subsection{Why White-Box Programs}

Source code analysis enables sound static flow checking before execution (black-box code cannot be reliably analyzed). Pre-execution rejection is stronger than runtime monitoring. Source checking enables compositional security: the server independently verifies program safety regardless of the author's trustworthiness.

\subsection{Declassification Design}

Declassification only via trusted server transformations (not user code) ensures explicit target policies, authority proof from owners, transform-specific constraints, and immutable audit logs. This centralizes policy enforcement and prevents user programs from bypassing controls.

\subsection{Attack-Driven Justification}

We design the system to withstand realistic attacks. Below we enumerate attack classes and show where each is blocked:

\textbf{A1: Direct Copy Leak.} An attacker with authorized access to high-confidentiality data $secret: (owner:\{owner\})$ attempts to copy it to low-confidentiality output $public: (owner:\{\})$ via assignment $public := secret$.

The type checker must verify: $(owner:\{owner\}) \sqsubseteq (owner:\{\})$? This requires $\{owner\} \supseteq \{\}$, which is true. So the flow would be accepted!

Wait—this seems wrong. The issue is that low confidentiality should be on the right-hand side of the assignment check. Actually, in our type system, we require $L_e \sqsubseteq \Gamma(x)$, so $L_{secret} \sqsubseteq L_{public}$. For this to hold, $(owner:\{owner\}) \sqsubseteq (owner:\{\})$, which fails because $\{owner\} \not\supseteq \{\}$. The assignment is blocked. \checkmark

\textbf{A2: Implicit Control Leak via Branching.} An attacker crafts a program:
\[\text{if } (secret > threshold) \text{ then } result := 1 \text{ else } result := 0\]

where $secret$ is high-confidentiality. An observer watching $result$ could infer something about $secret$ by seeing whether $result$ is 1 or 0.

The checker raises $pc$ inside the branch to $pc \sqcup L_{secret}$. Any assignment to $result$ inside the branch must then satisfy $(pc \sqcup L_{secret}) \sqsubseteq L_{result}$. If $L_{secret}$ is high-confidentiality and $L_{result}$ is low-confidentiality, this check fails. The program is rejected. \checkmark

\textbf{A3: Integrity Poisoning.} A low-integrity attacker (researcher with limited writing authority) attempts to corrupt high-integrity data by supplying it as input to a program whose output is then written to high-integrity storage.

The flow rule prevents this: if the output requires high integrity, $L_{req}$ will have integrity components with few trusted writers. The program's input from the low-trust researcher will not satisfy $L_{prog\_input} \sqsubseteq L_{req}$ in the integrity component, because the researcher is not in the high-integrity writer set. \checkmark

\textbf{A4: Unauthorized Downgrade.} An attacker without authority claims to declassify sensitive data from $(patient:\{patient,doctor\})$ to $(patient:\{patient,doctor,researcher\})$.

The server checks: ``does the caller have authority from the patient to add researcher to the reader set?'' If not, the call fails. Declaring authority with an invalid credential is rejected. \checkmark

\subsection{Boundaries of the Argument}

Our proof obligations do not cover:
\begin{itemize}
    \item compromised or malicious trusted server code,
    \item covert channels outside the language and API model (e.g., timing attacks, power analysis),
    \item social engineering or misrepresentation of authority by legitimate users.
\end{itemize}

Explicitly stating these boundaries keeps our claims accurate and focused on what the model and implementation demonstrably provide.

